\Footer{Zhuohang Bian · Curriculum Vitae}{\thepage}

\Header{Zhuohang Bian}
  {Ph.D. Student, Peking University · LLM Inference Systems and Computer Architecture}
  {%
   \Contact{\href{mailto:bianzhuohang26@stu.pku.edu.cn}{bianzhuohang26@stu.pku.edu.cn}}%
   \Contact{+86 136 5616 0098}%
   \Contact{\href{https://github.com/zhhangBian}{github.com/zhhangBian}}%
   \Contact{\href{https://zhhangbian.github.io/}{zhhangbian.github.io}}}

\Section{Education}
\Entry{Peking University · School of Integrated Circuits}{2026 – Present}
\Sub{Ph.D. in Electronic Design Automation and Computing Systems}
\Entry{Beihang University · School of Computer Science and Engineering}{2022 – 2026}
\Sub{B.Eng. in Computer Science and Technology · GPA 3.86/4.00 · Rank 19/239 (top 10\%)}

\Section{Publications}
\Pub{TokenCake: A KV-Cache-centric Serving Framework for LLM-based Multi-Agent Applications}
  {\Me{Zhuohang Bian}, Feiyang Wu, Zhuoran Li, Teng Ma, Youwei Zhuo}
  {EuroSys 2027}
  {\href{https://arxiv.org/abs/2510.18586}{[Paper]}}
\Pub{Graph.hls: A Compiler Framework for Composable Graph Accelerator Design}
  {Feiyang Wu, Xuxiao Yang, \Me{Zhuohang Bian}, Jing Wang, Ruifan Xu, Guangyu Sun, Yun Liang, Youwei Zhuo}
  {ISCA 2026}
  {\href{https://doi.org/10.1109/ISCA66397.2026.00064}{[Paper]}}
\Pub{MammoExpert: Benchmarking Chain-of-Thought Reasoning in Mammography Diagnosis}
  {Di Dai, Bo Liu, Youcheng Li, Haojun Yu, \Me{Zhuohang Bian}, et al.}
  {KDD 2026}
  {\href{https://arxiv.org/abs/2606.21119}{[Paper]}}
\Pub{Empirical Study on Robustness and Resilience in Cooperative Multi-Agent Reinforcement Learning}
  {Simin Li, Zihao Mao, Hanxiao Li, Zonglei Jing, \Me{Zhuohang Bian}, et al.}
  {NeurIPS 2025}
  {\href{https://arxiv.org/abs/2510.11824}{[Paper]}}
\Pub{TokenSim: Enabling Hardware and Software Exploration for Large Language Model Inference Systems}
  {Feiyang Wu, \Me{Zhuohang Bian}, Guoyang Duan, Tianle Xu, Junchi Wu, Teng Ma, Yongqiang Yao, Ruihao Gong, Youwei Zhuo}
  {APPT 2025}
  {\href{https://arxiv.org/abs/2503.08415}{[Paper]}}
\Pub{TokenDance: Scaling Multi-Agent LLM Serving via Collective KV Cache Sharing}
  {\Me{Zhuohang Bian}, Feiyang Wu, Chengrui Zhang, Hangcheng Dong, Yun Liang, Youwei Zhuo}
  {Preprint}
  {\href{https://arxiv.org/abs/2604.03143}{[Paper]}}
\Pub{TokenStack: A Heterogeneous HBM-PIM Architecture and Runtime for Efficient LLM Inference}
  {Zhuoran Li, \Me{Zhuohang Bian}, Zihao Huang, Yibo Zhao, Xueqi Li, Guangyu Sun, Youwei Zhuo}
  {Preprint}
  {\href{https://arxiv.org/abs/2605.05639}{[Paper]}}
\Pub{A Full-Stack Characterization of High-Bandwidth Flash for KV-Centric LLM Serving}
  {Zhuoran Li, \Me{Zhuohang Bian}, Xin Huang, Yibo Zhao, Guangyu Sun, Youwei Zhuo}
  {Preprint}
  {\href{https://arxiv.org/abs/2608.11668}{[Paper]}}

\Section{Research Experience}
\Entry{Peking University · Prof.~Youwei Zhuo's Group}{Feb 2025 – Present}
\Sub{Multi-agent LLM inference systems: KV-cache scheduling, reuse and compression}

\Entry{TokenCake: Agent-Aware KV-Cache Scheduling}{First author · EuroSys 2027}
\Item{Modeled workflow dependencies and tool calls to expose KV caches that sit idle during tool execution and contention on the critical path; designed \textbf{coordinated temporal and spatial schedulers}.}
\Item{Combined proactive offloading, predictive reloading and dynamic GPU memory partitioning, reducing end-to-end latency by \textbf{up to 47.06\%} versus vLLM in memory-constrained multi-agent workloads.}
\Item{Raised GPU KV-cache utilization by up to \textbf{16.9 percentage points} over vLLM (Qwen2.5-14B, Code-Writer). \Links{\href{https://arxiv.org/abs/2510.18586}{[Paper]} \href{https://github.com/pku-lemonade/TokenCake}{[Code]}}}

\Entry{TokenDance: Collective KV-Cache Reuse and Compression}{First author · Under review}
\Item{Designed a \textbf{KV Collector} for round-level reuse of shared context across agents, removing the position shifts that defeat prefix caching and cause redundant prefill.}
\Item{Designed \textbf{Diff-Aware Storage}, encoding similar KV caches as block-sparse differences for \textbf{11–17× compression}.}
\Item{Served up to \textbf{2.7× more concurrent agents} than vLLM with prefix caching under the latency SLO, with up to \textbf{1.9× faster prefill} than per-request position-independent caching. \Links{\href{https://arxiv.org/abs/2604.03143}{[Paper]}}}

\Entry{TokenSim: LLM Inference System Simulator}{Second author · APPT 2025}
\Item{Co-developed a modular simulator with extensible schedulers, GPU memory managers and operator-level performance models for design-space exploration of P/D disaggregation and heterogeneous hardware.}
\Item{Matched real vLLM runs with \textbf{under 1\% geometric-mean error} in throughput and P50/P99/max latency (A100, LLaMA2-7B, ShareGPT). \Links{\href{https://arxiv.org/abs/2503.08415}{[Paper]} \href{https://github.com/pku-lemonade/TokenSim}{[Code]}}}

\Entry{Beihang University · Prof.~Xianglong Liu's Group}{Jul 2024 – Jan 2025}
\Sub{Robustness of cooperative multi-agent reinforcement learning}
\Item{Implemented attack algorithms and experiment pipelines for a benchmark of \textbf{82,620 experiments} across 4 environments, 13 uncertainty types and 15 key hyperparameters; contributed to analysis and writing (NeurIPS 2025).}

\Entry{Beihang University · Prof.~Zhongzhi Luan's Group}{Sep 2023 – May 2024}
\Sub{Transformer operators on the MT-3000 CPU-DSP heterogeneous processor}
\Item{Implemented and optimized the Linear operator on MT-3000 DSPs with matrix tiling and DMA double buffering, integrated through PyTorch C++ extensions.}
\Item{Contributed to multi-head attention integration; the team's implementation reached up to \textbf{34× speedup} over a general-purpose CPU core. Outstanding Paper Award, HPC China 2024.}

\Section{Industry Experience}
\Entry{SenseTime · LLM Inference Systems Intern, \href{https://github.com/ModelTC/LightLLM}{LightLLM}}{Jun 2025 – Oct 2025}
\Sub{Contributed 6 upstream pull requests: P/D-disaggregated serving, KV-cache transfer, multimodal models, observability}
\Item{Refactored prefill/decode (P/D) node selection into extensible selectors and added \textbf{memory-aware routing} driven by predicted node load, with load reporting from each node. \href{https://github.com/ModelTC/LightLLM/pull/970}{\#970}}
\Item{Implemented \textbf{timeout retries and KV-cache return transfers} for chunked decoding, and \textbf{CUDA/NCCL-based cross-GPU KV-cache migration} across multiple data-parallel ranks. \href{https://github.com/ModelTC/LightLLM/pull/944}{\#944} \href{https://github.com/ModelTC/LightLLM/pull/999}{\#999}}
\Item{Integrated the \textbf{MinerU2} multimodal model end to end: tokenizer, vision modules, vision-service RPC and serving APIs. \href{https://github.com/ModelTC/LightLLM/pull/1034}{\#1034}}
\Item{Added multimodal token-usage accounting and service-process identifiers for service observability. \href{https://github.com/ModelTC/LightLLM/pull/1011}{\#1011} \href{https://github.com/ModelTC/LightLLM/pull/1024}{\#1024}}

\Section{Projects}
\Entry{BOOMCore: Out-of-Order LoongArch Processor}{Jul 2024 – Aug 2024}
\Item{Co-developed a \textbf{13-stage dual-issue superscalar processor} (LoongArch32R) with branch prediction, L1 cache and TLB that boots Linux at \textbf{80 MHz} on FPGA; 4th nationally in the Loongson Cup.}
\Item{Owned the ALU/LSU/MDU issue queues, commit logic and memory-access state machines, and backend integration. \Links{\href{https://github.com/zhhangBian/BOOM_chip}{[Code]}}}
\Entry{Compiler Course Infrastructure}{Beihang University}
\Item{Built a complete compiler course project that was adopted as a teaching example in national compiler-course teacher training. \Links{\href{https://github.com/zhhangBian/Compiler-Techniques}{[Code]}}}

\Section{Skills}
\Skill{Languages \& GPU}{Python, C++, CUDA, NCCL, PyTorch C++ extensions}
\Skill{Inference systems}{LightLLM, vLLM, SGLang; KV-cache management, P/D disaggregation, request scheduling, cross-GPU migration, profiling}
\Skill{Architecture}{SystemVerilog, FPGA, out-of-order processors, CPU-DSP heterogeneous computing, hardware/software performance modeling}
\Skill{English}{CET-6 (589)}

\Section{Honors \& Service}
\Award{Outstanding Graduate of Beijing}{2026}
\Award{Tutorial Chair, APPT (T2: Graph.hls)}{2026}
\Award{National Second Prize, National Student Computer System Capability Challenge}{2025}
\Award{National First Prize, China Undergraduate Mathematical Contest in Modeling}{2024}
\Award{Meritorious Winner, Mathematical Contest in Modeling (MCM/ICM)}{2024}
\Award{4th Place Nationally (Second Prize), Loongson Cup CPU Track}{2024}
\Award{Teaching Assistant, Computer Organization and Operating Systems, Beihang University}{2024 – 2025}
